% concatenated from Yang_MatrixKloostermanSum_2025.tex
\documentclass[11pt]{article}
\usepackage{geometry}
\geometry{a4paper, margin=1in}
\usepackage{amsmath, amssymb, amsfonts, amsthm}
\usepackage{graphicx}
\usepackage{hyperref}
\usepackage{booktabs}
\usepackage{bm}
\usepackage{parskip}
\usepackage[ruled,vlined]{algorithm2e}
\usepackage{youngtab}
\usepackage{ytableau}
\usepackage{enumerate}
\usepackage{mathrsfs}

% Theorem environments
\newtheorem{theorem}{Theorem}[section]
\newtheorem{lemma}[theorem]{Lemma}
\newtheorem{proposition}[theorem]{Proposition}
\newtheorem{corollary}[theorem]{Corollary}
\newtheorem{definition}[theorem]{Definition}
\newtheorem{example}[theorem]{Example}
\newtheorem{remark}[theorem]{Remark}
\newtheorem{conjecture}[theorem]{Conjecture}  % Added conjecture environment

% Custom commands
\newcommand{\Fq}{\mathbb{F}_q}
\newcommand{\Fqd}{\mathbb{F}_{q^d}}
\newcommand{\GL}{\mathrm{GL}}
\newcommand{\Kl}{\mathrm{Kl}}
\newcommand{\tKl}{\widetilde{\mathrm{Kl}}}
\newcommand{\Sym}{\mathrm{Sym}}
\newcommand{\Tr}{\mathrm{Tr}}
\newcommand{\Det}{\mathrm{Det}}
\newcommand{\Char}{\mathrm{Char}}
\newcommand{\Fr}{\mathrm{Fr}}
\newcommand{\diag}{\mathrm{diag}}
\newcommand{\Irr}{\mathrm{Irr}}
\newcommand{\cusp}{\mathrm{cusp}}

% Group theory commands
\newcommand{\SL}{\mathrm{SL}}
\newcommand{\SU}{\mathrm{SU}}
\newcommand{\Sp}{\mathrm{Sp}}
\newcommand{\U}{\mathrm{U}}

\title{Matrix Kloosterman Sums, Random Matrix Statistics, and Cryptography}
\author{Tianshuo Yang}
\date{}

\begin{document}

\maketitle

\begin{abstract}
This paper presents a comprehensive study of matrix Kloosterman sums, including their computational aspects, distributional behavior, and applications in cryptographic analysis. Building on the work of \cite{zelingher2023matrix}, we develop algorithms for evaluating these sums via Green’s polynomials and establish a general framework for analyzing their statistical distributions. We further investigate the associated $L$-functions and clarify their relationships with symmetric functions and random matrix theory.
We show that, analogous to the eigenvalue statistics of random matrices in compact Lie groups such as $SU(n)$ and $Sp(2n)$, the normalized values of matrix Kloosterman sums exhibit Sato–Tate equidistribution. 
%This phenomenon is formalized into a statistical testing framework. 
Finally, we apply this framework to distinguish truly random sequences from those exhibiting subtle algebraic biases, and we propose a novel spectral test for cryptographic security based on the distributional signatures of matrix Kloosterman sums.
\end{abstract}

\section{Introduction}

Kloosterman sums are fundamental objects in number theory, algebraic geometry, and representation theory. Classically, they are defined as exponential sums of the form
\[
K(a,b;p) = \sum_{\substack{x \in \mathbb{F}_p^\times \\ xy = 1}} e^{2\pi i (ax + by)/p},
\]
and they exhibit exhibit deep connections with modular forms, automorphic representations, and $\ell$-adic cohomology \cite{Katz88}. Their natural extension to matrix groups—known as matrix Kloosterman sums—has been investigated in recent work of Erd\'{e}lyi and T\'{o}th \cite{erdelyi2021matrix}, and independently by Zelingher \cite{zelingher2023matrix}. These developments have opened new avenues in the study of random matrix statistics, equidistribution phenomena, and applications to cryptographic analysis.

Matrix Kloosterman sums are defined for matrices in the general linear group $\GL_n(\Fq)$ and incorporate both multiplicative characters (through determinants) and additive characters (through traces). Their significance stems from several perspectives:

\begin{enumerate}
\item \textbf{Algebraic Structure}: Matrix Kloosterman sums are intimately connected with the representation theory of $\GL_n(\Fq)$ and symmetric functions through Green polynomials (\cite{green1955characters, macdonald1998Symmetric}, and \cite{zelingher2023matrix}).

\item \textbf{Random Matrix Statistics}: Deligne's equidistribution theorem \cite{Katz88} suggests that normalized Kloosterman sums should be distributed like traces of random matrices in compact Lie groups. This connection between number theory and random matrix theory has profound implications for understanding the statistical properties of these sums.

\item \textbf{Cryptographic Applications}: The randomness properties of matrix Kloosterman sums make them suitable for testing cryptographic algorithms. By comparing the distribution of matrix Kloosterman sums computed from cipher outputs with theoretical random matrix distributions, one can assess the quality of cryptographic randomness.
\end{enumerate}

In this paper, we develop algorithms for computing matrix Kloosterman sums using Green polynomials and classical Kloosterman sums, building on the foundational work of Zelingher \cite{zelingher2023matrix}. We design and implement two complementary algorithms: Algorithm I, which computes matrix Kloosterman sums for Jordan-type matrices, and Algorithm II, which extends the computation to arbitrary matrices via their Jordan decomposition.

We investigate the $L$-functions associated with matrix Kloosterman sums and establish their connections with symmetric functions and random matrix statistics. Through extensive numerical experiments, we confirm the Sato–Tate distribution behavior of normalized matrix Kloosterman sums and demonstrate their equidistribution in compact Lie groups. Using these insights, we propose a novel statistical test for cryptographic randomness based on the distribution of matrix Kloosterman sums and show its effectiveness in detecting algebraic or structural deviations from true randomness.

The paper is organized as follows. Section~2 reviews the necessary background on conjugacy classes and representations of $\GL_n(\Fq)$, symmetric functions, classical Kloosterman sums, and matrix Kloosterman sums. Section~3 presents the computational algorithms. Section~4 investigates the associated $L$-functions. Section~5 explores the connections with random matrix theory. Section~6 describes cryptographic applications and numerical experiments. Section~7 concludes with directions for future research.


\section{Preliminaries}
\label{sec:preliminaries}

\subsection{Conjugacy Classes of $\GL_n(\Fq)$}

Let $\Fq$ be a finite field of size $q$, where $q$ is a power of some prime number $p$. The classification of conjugacy classes in $\GL_n(\Fq)$ is fundamental to our study. We follow Green's approach \cite{green1955characters}.
\medskip

\begin{definition}[Regular Elliptic Element]
An element $g \in \GL_n(\Fq)$ is called \emph{regular elliptic} if its characteristic polynomial is irreducible over $\Fq$. Such an element has $n$ distinct eigenvalues in $\mathbb{F}_{q^n}$: $\{\xi, \xi^q, \ldots, \xi^{q^{n-1}}\}$.
\end{definition}
\medskip

\begin{theorem}[Conjugacy Classification] \cite{green1955characters}
For any $y \in \GL_n(\Fq)$, there exists an integer $s \geq 1$, $d_1, \ldots, d_s \geq 1$, regular elliptic elements $y_j \in \GL_{d_j}(\Fq)$ (with mutually disjoint eigenvalues over $\overline{\Fq}$), and non-empty partitions $\bm{\mu}_1, \ldots, \bm{\mu}_s$, such that
\[
y \sim \diag\left(J_{\bm{\mu}_1}(y_1), \ldots, J_{\bm{\mu}_s}(y_s)\right),
\]
where $J_{\bm{\mu}_j}(y_j)$ denotes the Jordan block matrix associated with partition $\bm{\mu}_j$ and regular elliptic element $y_j$.
\end{theorem}
\medskip

The Jordan block $J_{(\mu)}(y)$ for a single part $\mu$ and regular elliptic element $y \in \GL_d(\Fq)$ is defined as:
\[
J_{(\mu)}(y) = \begin{pmatrix}
	y & I_d & & \\
	& y & \ddots & \\
	& & \ddots & I_d \\
	& & & y
\end{pmatrix} \in \GL_{d\mu}(\Fq),
\]
where $I_d$ is the $d \times d$ identity matrix.

\subsection{Representations of $\GL_n(\Fq)$}

The representation theory of $\GL_n(\Fq)$ was completely described by Green \cite{green1955characters}. Let $\Irr(\GL_n(\Fq))$ denote the set of irreducible representations and $\Irr_{\cusp}(\GL_d(\Fq))$ the subset of cuspidal representations.
\medskip

\begin{theorem}[Green's Parameterization] \cite{green1955characters}
There is a bijection between $\Irr(\GL_n(\Fq))$ and parameters $\varphi$ with $|\varphi| = n$, where a parameter $\varphi$ assigns to each pair $(d,\sigma)$ with $\sigma \in \Irr_{\cusp}(\GL_d(\Fq))$ a partition $\varphi(d,\sigma)$, such that
\[
|\varphi| = \sum_{(d,\sigma)} d \cdot |\varphi(d,\sigma)|.
\]
\end{theorem}

Cuspidal representations correspond to Frobenius orbits of regular characters $\theta: \Fqd^\times \to \mathbb{C}^\times$. A character $\theta$ is \emph{regular} if $\theta^{q^j} \neq \theta$ for all $0 \leq j \leq d-1$, meaning its Frobenius orbit has size $d$.

\subsection{Symmetric Functions and Green Polynomials}

We recall essential facts about symmetric functions and Green polynomials \cite{macdonald1998Symmetric}.
\medskip

\begin{definition}[Partitions]
A partition $\bm{\lambda} = (\lambda_1, \ldots, \lambda_\ell)$ is a weakly decreasing sequence of positive integers. We write $\bm{\lambda} \vdash r$ if $|\bm{\lambda}| = \sum \lambda_i = r$. The length $\ell(\bm{\lambda}) = \ell$, and $m_i(\bm{\lambda})$ counts parts equal to $i$.
\end{definition}

For a partition $\bm{\lambda}$, define $z_{\bm{\lambda}} = \prod_{i \geq 1} i^{m_i} m_i!$ and $n(\bm{\lambda}) = \sum_{i=1}^{\ell} (i-1)\lambda_i$. 
%\begin{align*}
%z_{\bm{\lambda}} &= \prod_{i \geq 1} i^{m_i} m_i!,\\
%\epsilon_{\bm{\lambda}} &= (-1)^{|\bm{\lambda}| - \ell(\bm{\lambda})},\\
%n(\bm{\lambda}) &= \sum_{i=1}^{\ell} (i-1)\lambda_i.
%\end{align*}
We use $p_{\bm{\lambda}}$ to denote the power-sum symmetric function associated to the partition $\bm{\lambda}$ and $P_{\bm{\mu}}(x;t)$ to denote the Hall–Littlewood polynomial associated to the partition $\bm{\mu}$.
\medskip

\begin{definition}[Green Polynomials]
The Green polynomials $Q_{\bm{\lambda}}^{\bm{\mu}}(q)$ are defined by:
\[
Q_{\bm{\lambda}}^{\bm{\mu}}(q) = q^{n(\bm{\mu})} X_{\bm{\lambda}}^{\bm{\mu}}(q^{-1}),
\]
where $X_{\bm{\lambda}}^{\bm{\mu}}(t)$ are determined by the expansion:
\[
p_{\bm{\lambda}}(x) = \sum_{\bm{\mu}} X_{\bm{\lambda}}^{\bm{\mu}}(t) P_{\bm{\mu}}(x; t).
\]
\end{definition}

Green polynomials satisfy important properties \cite{green1955characters}:
\begin{itemize}
\item $Q_{\bm{\lambda}}^{\bm{\mu}}(t) \in \mathbb{Z}[t]$ and vanish unless $|\bm{\lambda}| = |\bm{\mu}|$.
\item $Q_{\bm{\lambda}}^{(r)}(t) = 1$ for all partitions $\bm{\lambda}$ of $r$.
\item $\sum_{|\bm{\lambda}|=r} z_{\bm{\lambda}}^{-1} Q_{\bm{\lambda}}^{\bm{\mu}}(q) = 1$.
\item ${\rm deg}\left(Q_{\bm{\lambda}}^{\bm{\mu}}(t)\right)\leq n(\bm{\mu})=\sum_{i=1}^{\ell}(i-1)\mu_i$ 
\end{itemize}
For tables of Green polynomials, see Green \cite{green1955characters} for $n\leq 5$ and Morris \cite{morris1963characters} for $n=6$ and $7$.

\subsection{Classical Kloosterman Sums and Sheaves}

Let $\psi: \Fqd \to \mathbb{C}^\times$ be a nontrivial additive character and $\bm{\alpha} = (\alpha_1, \ldots, \alpha_k)$ with $\alpha_i: \Fqd^\times \to \mathbb{C}^\times$ multiplicative characters.
\medskip

\begin{definition}[Hyper-Kloosterman Sum]
The rank $k$ Kloosterman sum is:
\[
\Kl(\bm{\alpha}, \psi, \xi) = (-1)^{k-1} \sum_{\substack{t_1, \ldots, t_k \in \Fqd^\times \\ t_1 \cdots t_k = \xi}} \alpha_1(t_1) \cdots \alpha_k(t_k) \psi(t_1 + \cdots + t_k).
\]
\end{definition}

Katz \cite{Katz88} studied Kloosterman sheaves, which are $\ell$-adic sheaves on the multiplicative group $\mathbb{G}_m$ whose Frobenius traces give Kloosterman sums. For a rank $k$ Kloosterman sheaf $\Kl_k^{\Fqd}(\bm{\alpha},\psi)$, the Frobenius eigenvalues at $\xi \in \Fqd^\times$ are denoted $w_1, \ldots, w_k$, satisfying $|w_j| = q^{d(k-1)/2}$.

\subsection{Matrix Kloosterman Sums}
\label{sec:matrix_kl}

We recall the definition of matrix Kloosterman sums and their properties from \cite{erdelyi2021matrix} and \cite{zelingher2023matrix}.

\subsubsection{Definition}

\begin{definition}[Matrix Kloosterman Sums]
Let $y \in \GL_n(\Fq)$, $\psi: \Fq \to \mathbb{C}^\times$ a nontrivial additive character, and $\bm{\alpha} = (\alpha_1, \ldots, \alpha_k)$ multiplicative characters of $\Fq^\times$. The matrix Kloosterman sum is defined as:
\[
\Kl_n(\bm{\alpha}, \psi, y) = \sum_{\substack{g_1, \ldots, g_k \in \GL_n(\Fq) \\ \prod_{j=1}^k g_j = y}} \left( \prod_{i=1}^k \alpha_i(\Det(g_i)) \right) \psi\left( \Tr\left( \sum_{i=1}^k g_i \right) \right).
\]
\end{definition}

This definition generalizes classical Kloosterman sums in two directions: from scalars to matrices, and from products to both products (through the condition $\prod g_j = y$) and sums (through the trace in the additive character).

\subsubsection{Reduction Formulas}

Matrix Kloosterman sums exhibit multiplicative properties under direct sums of matrices with disjoint eigenvalues.
\medskip

\begin{theorem}[Direct Sum Decomposition] \cite{erdelyi2021matrix}
\label{thm:direct_sum}
Let $y_1 \in \GL_{n_1}(\Fq)$ and $y_2 \in \GL_{n_2}(\Fq)$ have no common eigenvalues over $\overline{\Fq}$. Then:
\[
\Kl_{n_1+n_2}(\bm{\alpha}, \psi, \diag(y_1, y_2)) = q^{(k-1)n_1n_2} \Kl_{n_1}(\bm{\alpha}, \psi, y_1) \Kl_{n_2}(\bm{\alpha}, \psi, y_2).
\]
\end{theorem}

For regular elliptic elements, matrix Kloosterman sums reduce to classical Kloosterman sums.
\medskip

\begin{theorem}[Regular Elliptic Reduction]\cite{zelingher2023matrix}
\label{thm:reg_elliptic}
Let $y \in \GL_n(\Fq)$ be regular elliptic with eigenvalues $\{\xi, \xi^q, \ldots, \xi^{q^{n-1}}\}$, $\xi \in \mathbb{F}_{q^n}^\times$. Then:
\[
\Kl_n(\bm{\alpha}, \psi, y) = (-1)^{(k-1)(n+1)} q^{(k-1)\binom{d}{2}} \Kl_n(\bm{\alpha}, \psi, \xi),
\]
where the right side is the classical Kloosterman sum in $\mathbb{F}_{q^n}$.
\end{theorem}

\subsubsection{Jordan Block Formulas}

The most intricate case involves Jordan blocks. Let $n = dr$ and $y \in \GL_d(\Fq)$ be regular elliptic with eigenvalues $\{\xi, \xi^q, \ldots, \xi^{q^{d-1}}\}$.
\medskip

\begin{theorem}[Jordan Block Formula] \cite{zelingher2023matrix}
\label{thm:jordan_formula}
For any partition $\bm{\mu} \vdash r$:
\[
\Kl_n(\bm{\alpha}, \psi, J_{\bm{\mu}}(y)) = (-1)^{(k-1)n} q^{(k-1)\binom{n}{2}} \sum_{\bm{\lambda} \vdash r} \frac{Q_{\bm{\lambda}}^{\bm{\mu}}(q^d)}{z_{\bm{\lambda}}} (-1)^{(k-1)\ell(\bm{\lambda})} \prod_{j=1}^{\ell(\bm{\lambda})} \Kl_{d\lambda_j}(\bm{\alpha}, \psi, \xi).
\]
\end{theorem}

Let $w_1, \ldots, w_k$ be the Frobenius eigenvalues of $\Kl_k^{\Fqd}(\bm{\alpha},\psi)$ at $\xi$, and define normalized eigenvalues $\widetilde{w}_j = q^{-d(k-1)/2} w_j$. Using the identity
\[
p_{\bm{\lambda}}(w_1, \ldots, w_k) = (-1)^{(k-1)\ell(\bm{\lambda})} \prod_{j=1}^{\ell(\bm{\lambda})} \Kl_{d\lambda_j}(\bm{\alpha}, \psi, \xi),
\]
One obtains a symmetric function formulation:
\medskip

\begin{corollary}
\[
\Kl_n(\bm{\alpha}, \psi, J_{\bm{\mu}}(y)) = (-1)^{(k-1)n} q^{(k-1)\binom{n}{2}} \sum_{\bm{\lambda} \vdash r} \frac{Q_{\bm{\lambda}}^{\bm{\mu}}(q^d)}{z_{\bm{\lambda}}} p_{\bm{\lambda}}(w_1, \ldots, w_k).
\]
\end{corollary}

This reveals that matrix Kloosterman sums are essentially evaluations of symmetric functions at the Frobenius eigenvalues of Kloosterman sheaves, with Green polynomials as coefficients.

\subsubsection{General Formula}

Combining Theorems \ref{thm:direct_sum} and \ref{thm:jordan_formula}, one obtains a complete formula for arbitrary matrices.
\medskip

\begin{theorem}[General Formula] \cite{zelingher2023matrix}
\label{thm:general_formula}
Let $y = \diag(J_{\bm{\mu}_1}(y_1), \ldots, J_{\bm{\mu}_s}(y_s))$, where $y_j \in \GL_{d_j}(\Fq)$ are regular elliptic with mutually disjoint eigenvalues, and $\bm{\mu}_j \vdash r_j$ with $\sum_{j=1}^s d_j r_j = n$. Let $\eta_j \in \mathbb{F}_{q^{d_j}}^\times$ correspond to eigenvalues of $y_j$, and $w_{j1}, \ldots, w_{jk}$ be Frobenius eigenvalues of $\Kl_k^{\mathbb{F}_{q^{d_j}}}(\bm{\alpha},\psi)$ at $\eta_j$. Then:
$$
\Kl_n(\bm{\alpha}, \psi, y) = (-1)^{(k-1)n} q^{(k-1)\binom{n}{2}} \times
\prod_{j=1}^s \sum_{\bm{\lambda} \vdash r_j} \frac{Q_{\bm{\lambda}}^{\bm{\mu}_j}(q^{d_j})}{z_{\bm{\lambda}}} p_{\bm{\lambda}}(w_{j1}, \ldots, w_{jk}).
$$
\end{theorem}

In the special case $\bm{\mu} = (r)$, i.e., $J_{\bm{\mu}}(y)$ is a regular element, we have that $Q^{(r)}_{\bm{\lambda}}(t)=1$ for all $\bm{\lambda}\vdash r$ and therefore 
\[
\sum_{\bm{\lambda} = (\lambda_1, \ldots, \lambda_{\ell}) \vdash (r)}\frac{Q_{\bm{\lambda}}^{(r)}(t)}{z_{\bm{\lambda}}}p_{\bm{\lambda}}(x_1, \ldots, x_k)= h_r(x_1, \dots, x_k)=\sum_{\bm{\lambda} \vdash r} m_{\bm{\lambda}}(x_1, \dots, x_k)
\]
is the complete symmetric polynomial of degree $r$.

\section{Algorithms for Computation}
\label{sec:algorithms}

\subsection{Normalization}

For statistical analysis, we work with normalized matrix Kloosterman sums:
\[
\tKl_n(\bm{\alpha}, \psi, y) = q^{-(k-1)\frac{n^2}{2}} \Kl_n(\bm{\alpha}, \psi, y).
\]
By Green \cite{green1955characters}, ${\rm deg}\left(Q_{\bm{\lambda}}^{\bm{\mu}}(t)\right)\leq n(\bm{\mu})=\sum_{i=1}^{\ell}(i-1)\mu_i$. For all partition $\bm{\mu}\vdash r$, the maximum $n(\bm{\mu})$ is $\binom{r}{2}$, achieved by $\bm{\mu}=(1, 1, \ldots, 1)$. We have that $|\tKl_n(\bm{\alpha}, \psi, y)| \leq q^{\sum_{j=1}^s\frac{d_jr_j(r_j-1)}{2}}$. In particular, for $\bm{\mu}=(r)$, we have $Q_{\bm{\lambda}}^{(r)}(t) = 1$ for all partitions $\bm{\lambda}$ of $r$. Hence for Jordan block $J_{(r)}(y)$, where $y$ is a regular elliptic element, the normalized sums are bounded independently of $q$.

\subsection{Algorithm I: Jordan Block Computation}

\begin{algorithm}[H]
\caption{Compute $\tKl_n(\bm{\alpha}, \psi, J_{\bm{\mu}}(y))$}
\label{alg:jordan}
\KwIn{$q = p^m$, $n$, regular elliptic $y \in \GL_d(\Fq)$, partition $\bm{\mu} \vdash r$ with $n = dr$}
\KwOut{$\tKl_n(\bm{\alpha}, \psi, J_{\bm{\mu}}(y))$}
\Begin{
Compute $P(T) = \det(TI_d - y)$ in $\Fq[T]$\;
Find $\xi \in \Fqd^\times$ such that $P(\xi) = 0$ and $\Fqd = \Fq[\xi]$\;
Initialize $\text{sum} = 0$\;
\For{each partition $\bm{\lambda} = (\lambda_1, \ldots, \lambda_\ell) \vdash r$}{
Load $Q_{\bm{\lambda}}^{\bm{\mu}}(t)$ (precomputed table)\;
Compute $z_{\bm{\lambda}} = \prod_{i \geq 1} i^{m_i} m_i!$\;
Compute $\widetilde{\Kl}_{d\lambda_j}(\bm{\alpha}, \psi, \xi)$ for $j = 1, \ldots, \ell$\;
$\text{term} = \frac{Q_{\bm{\lambda}}^{\bm{\mu}}(q^d)}{z_{\bm{\lambda}}} (-1)^{(k-1)\ell} \prod_{j=1}^\ell \widetilde{\Kl}_{d\lambda_j}(\bm{\alpha}, \psi, \xi)$\;
$\text{sum} = \text{sum} + \text{term}$\;
}
\Return{$(-1)^{(k-1)n} \cdot \text{sum}$}
}
\end{algorithm}

The algorithm requires precomputed tables of Green polynomials $Q_{\bm{\lambda}}^{\bm{\mu}}(t)$ for partitions of size up to $r$. Classical Kloosterman sums $\widetilde{\Kl}_{d\lambda_j}(\bm{\alpha}, \psi, \xi)$ admit efficient computation through established methods. In particular, the Bruhat decomposition of $\GL_n$ provides explicit summation formulas, while character-sum techniques—based on Gauss sums and multiplicative character expansions—further simplify the evaluation. Together, these tools form the basis for modern algorithms used in the computation of Kloosterman-type sums.

\subsection{Algorithm II: General Matrix Computation}

\begin{algorithm}[H]
\caption{Compute $\tKl_n(\bm{\alpha}, \psi, y)$ for general $y$}
\label{alg:general}
\KwIn{$q = p^m$, $n$, $y \in \GL_n(\Fq)$}
\KwOut{$\tKl_n(\bm{\alpha}, \psi, y)$}
\Begin{
Compute Jordan form: $y = \diag(J_{\bm{\mu}_1}(y_1), \ldots, J_{\bm{\mu}_s}(y_s))$\;
Initialize $\text{result} = 1$\;
\For{$j = 1$ to $s$}{
Let $d_j = \dim(y_j)$, $r_j = |\bm{\mu}_j|$\;
Compute $\eta_j \in \mathbb{F}_{q^{d_j}}^\times$ from $y_j$ as in Algorithm \ref{alg:jordan}\;
$\text{blockVal} = \tKl_{d_j r_j}(\bm{\alpha}, \psi, J_{\bm{\mu}_j}(y_j))$ using Algorithm \ref{alg:jordan}\;
$\text{result} = \text{result} \times \text{blockVal}$\;
}
\Return{$\text{result}$}
}
\end{algorithm}

\subsection{Complexity Analysis}

The computational complexity depends on several factors:
\begin{itemize}
\item \textbf{Jordan Decomposition}: Over a finite field $\mathbb{F}_q$, the rational and Jordan canonical forms of an $n \times n$ matrix can be computed in polynomial time. In particular, standard algorithms (e.g., Frobenius form followed by primary decomposition) run in $O(n^3)$ field operations.

\item \textbf{Green Polynomials}: 
Green polynomials $Q_{\boldsymbol{\lambda}}^{\boldsymbol{\mu}}(t)$ for $\mathrm{GL}_n(\mathbb{F}_q)$
can be computed effectively via their interpretation as the transition
coefficients between the Hall--Littlewood and power-sum bases in the
ring of symmetric functions \cite{macdonald1998Symmetric}.
Recursive formulas, vertex-operator constructions, and computer algebra
implementations allow one, for fixed $n$, to compute an individual
$Q_{\boldsymbol{\lambda}}^{\boldsymbol{\mu}}(t)$ in time polynomial in the
size of the partition data.

However, in applications to matrix Kloosterman sums, one typically needs
all Green polynomials of a given size.  For
$y \sim \operatorname{diag}(J_{\boldsymbol{\mu}_1}(y_1),\ldots,
J_{\boldsymbol{\mu}_s}(y_s))$, computation of
$\widetilde{\operatorname{Kl}}_n(\boldsymbol{\alpha},\psi,y)$ requires the
values $Q_{\boldsymbol{\lambda}}^{\boldsymbol{\mu}_j}(t)$ for all
partitions $\boldsymbol{\lambda}\vdash |\boldsymbol{\mu}_j|$ and
$1\le j\le s$.  Since the partition number $p(n)$ grows exponentially by
the Hardy--Ramanujan asymptotic, obtaining the full set
$\{ Q_{\boldsymbol{\lambda}}^{\boldsymbol{\mu}}(t): |\boldsymbol{\lambda}|=|\boldsymbol{\mu}|\le n \}$
poses a significant computational challenge.  It is therefore desirable
to develop more efficient algorithms for computing all Green polynomials
of bounded size.

\item \textbf{Classical Kloosterman Sums}: 
The naive computation of $\operatorname{Kl}_d(\boldsymbol{\alpha},\psi,\xi)$
requires $O(q^{d/2})$ operations, but can be accelerated using
character-sum techniques and structural decompositions, yielding
substantial improvements in practice.
\end{itemize}
For practical applications with moderate $n$ and $q$, these algorithms are feasible.

\section{$L$-Functions and Symmetric Functions}
\label{sec:lfunctions}

\subsection{$L$-Functions of Classical Kloosterman Sums}

For a fixed $\xi \in \Fqd^\times$, consider the generating series of normalized classical Kloosterman sums:
\[
L(T) = \exp\left( \sum_{n=1}^\infty \tKl_{dn}(\bm{\alpha}, \psi, \xi) \frac{T^n}{n} \right).
\]

\begin{theorem}(See for example \cite{fu2005lfunctions})
\label{thm:L_function}
$L(T)^{(-1)^k}$ is a polynomial of degree $k$:
\[
L(T)^{(-1)^k} = \prod_{i=1}^k (1 - \widetilde{w}_i T),
\]
where $\widetilde{w}_i = q^{-d(k-1)/2} w_i$ are normalized Frobenius eigenvalues.
\end{theorem}

%\begin{proof}
%Using the identity $p_n(\widetilde{w}_1, \ldots, \widetilde{w}_k) = (-1)^{k-1} \tKl_{dn}(\bm{\alpha}, \psi, \xi)$, we have:
%\begin{align*}
%\log L(T)^{(-1)^k} &= (-1)^k \sum_{n=1}^\infty \tKl_{dn}(\bm{\alpha}, \psi, \xi) \frac{T^n}{n} \\
%&= -\sum_{n=1}^\infty p_n(\widetilde{w}_1, \ldots, \widetilde{w}_k) \frac{T^n}{n} \\
%&= \sum_{i=1}^k \log(1 - \widetilde{w}_i T).
%\end{align*}
%Exponentiating gives the result.
%\end{proof}

\subsection{Coefficient Formulas}

Expanding the product gives expressions in terms of elementary symmetric functions:
\[
L(T)^{(-1)^k} = \sum_{r=0}^k (-1)^r e_r(\widetilde{w}_1, \ldots, \widetilde{w}_k) T^r.
\]

Using Newton's identities, the coefficients can be expressed in terms of the Kloosterman sums $K_j = \tKl_{dj}(\bm{\alpha}, \psi, \xi)$.
\medskip

\begin{proposition}
The coefficient $C_r$ of $T^r$ in $L(T)^{(-1)^k}$ is:
\[
C_r = \sum_{\substack{s_1, \ldots, s_r \geq 0 \\ 1s_1 + \cdots + rs_r = r}} \frac{K_1^{s_1} \cdots K_r^{s_r} (-1)^{k(s_1 + \cdots + s_r)}}{s_1! \cdots s_r! 1^{s_1} \cdots r^{s_r}}.
\]
\end{proposition}

This can be computed efficiently using partition enumeration.

\subsection{Matrix Kloosterman Sums and Generating Functions}

We have the following result about the generating function of matrix Kloosterman sums for the partition $(n)$:
\medskip

\begin{theorem} For regular elliptic $y \in \GL_d(\Fq)$, one has
\label{thm:generating_function}
\[
\sum_{n=0}^\infty (-1)^{(k-1)dn} \tKl_{dn}(\bm{\alpha}, \psi, J_{(n)}(y)) T^{dn} = \frac{1}{\prod_{j=1}^k (1 - \widetilde{w}_j T^d)}.
\]
\end{theorem}

\begin{proof}
When $\bm{\mu} = (n)$, we have $Q_{\bm{\lambda}}^{(n)}(t) = 1$ for all $\bm{\lambda} \vdash n$. Thus:
\[
\tKl_{dn}(\bm{\alpha}, \psi, J_{(n)}(y)) = (-1)^{(k-1)dn} \sum_{\bm{\lambda} \vdash n} \frac{1}{z_{\bm{\lambda}}} p_{\bm{\lambda}}(\widetilde{w}_1, \ldots, \widetilde{w}_k) = (-1)^{(k-1)dn} h_n(\widetilde{w}_1, \ldots, \widetilde{w}_k),
\]
where $h_n$ is the complete symmetric function. The generating function for complete symmetric functions is $\sum_{n=0}^\infty h_n T^n = \prod_{j=1}^k (1 - \widetilde{w}_j T)^{-1}$.
\end{proof}

This establishes a direct link between matrix Kloosterman sums and symmetric functions.

\subsection{Connection with Symmetric Power $L$-Functions}

The generating function in Theorem \ref{thm:generating_function} is essentially the $L$-function of the symmetric powers of the Kloosterman sheaf. Recall that for a sheaf $\mathcal{F}$, the $L$-function of its $\ell$-th symmetric power is:
\[
L(\Sym^\ell(\mathcal{F}), T) = \prod_{x \in |X|} \prod_{i_1 + \cdots + i_k = \ell} (1 - \alpha_1(x)^{i_1} \cdots \alpha_k(x)^{i_k} T^{\deg x})^{-1}.
\]

For any $y \in \mathrm{GL}_d(\mathbb{F}_q)$, let 
$\pi_1(\xi), \ldots, \pi_k(\xi)$ be all the eigenvalues of the 
geometric Frobenius element ${\rm Fr}_\xi$ on ${\rm Kl}_k^{\mathbb{F}_{q^d}}(\bm{\alpha}, \psi)_\xi$. By Theorem 4.3, we have the following result:
\medskip

\begin{corollary}
The generating function of matrix Kloosterman sums for partitions $(n)$ is related to the $L$-function of the first symmetric powers of the Kloosterman sheaf:
\[
L(\Sym^1(\Kl_k^{\Fqd}(\bm{\alpha},\psi)), T)==\prod_{\xi\in\overline{\mathbb{F}_{q}}^{\times}}\prod_{j=1}^k\frac{1}{1-\widetilde{\pi_j(\xi)}T^{{\rm deg}(\xi)}}
$$
$$
=\prod_{[y] \text{ regular elliptic}}\sum_{n=0}^{\infty}\widetilde{\operatorname{Kl}}_{{\rm deg}(y)\cdot n}(\bm{\alpha}, \psi, J_{(n)}(y))(-1)^{(k-1){\rm deg}(y)\cdot n}T^{{\rm deg}(y)\cdot n},
\]
where the right side is interpreted as a generating function for symmetric power $L$-functions.
\end{corollary}

\section{Random Matrix Statistics}
\label{sec:random_matrix}

\subsection{Deligne's Equidistribution Theorem}

Deligne's work on the Weil conjectures, extended by Katz, provides the theoretical foundation for the random matrix behavior of Kloosterman sums.
\medskip

\begin{theorem}[Deligne–Katz Equidistribution]\cite{Katz88}
Let $\Kl_k^{\Fqd}(\bm{\alpha},\psi)$ be a Kloosterman sheaf of rank $k$ with geometric monodromy group $G_{\rm geom}$. As $q \to \infty$ (or as we vary parameters), the Frobenius conjugacy classes ${\rm Fr}_\xi$ at points $\xi \in \Fqd^\times$ become equidistributed in the space of conjugacy classes of the compact group $G_{\rm geom}$, with respect to Haar measure.
\end{theorem}

%Based on Katz's work, the group \(G(\mathbb{C})\) “depends’’ on three parameters 
%\((p,n,\xi)\). In fact, it is independent of \(\xi\), and very nearly independent of \(p\) as well.
%\begin{enumerate}
%\item For $n \ge 2$ even, $p$ arbitrary, $G(\mathbb{C}) = \mathrm{Sp}(n)$ 
%\item For $n \ge 3$ odd, and $p$ odd, $G(\mathbb{C}) = \mathrm{SL}(n)$
%\item For $n \ge 3$ odd, $n\neq 7$, and $p = 2$, $G(\mathbb{C}) = \mathrm{SO}(n)$
%
%For $n=7$ odd and $p = 2$, $G(\mathbb{C}) = G_2$
%\end{enumerate}

For classical Kloosterman sums (the case \(k=2\)), Katz proved that the geometric
monodromy group of the associated Kloosterman sheaf is \(\SL(2)\) (equivalently,
its compact form \(\SU(2)\)). For higher-rank Kloosterman sums (\(k>2\)), the
geometric monodromy groups are typically the symplectic groups \(\Sp(2k)\) (for
even \(k\)) or closely related classical groups, depending on the parity of \(k\)
and the characteristic. Katz’s classification shows that, except for a few small
exceptional cases, the monodromy is as large as allowed by the symmetries of the
sheaf.

\subsection{Sato–Tate Distributions}

For $k=2$, the normalized Kloosterman sums $\tKl_d(\bm{\alpha}, \psi, \xi)$ lie in $[-2,2]$. Writing $\tKl_d(\bm{\alpha}, \psi, \xi) = 2\cos\theta$ with $\theta \in [0,\pi]$, the Sato–Tate law states:
\medskip

\begin{theorem}[Sato–Tate Distribution] \cite{Katz88}
As $q \to \infty$, the angles $\theta$ for classical Kloosterman sums ($k=2$) are distributed according to the density:
\[
f_{ST}(\theta) = \frac{2}{\pi} \sin^2\theta, \quad \theta \in [0,\pi].
\]
Equivalently, the normalized sums $x = 2\cos\theta$ have density:
\[
f_{trace}(x) = \frac{1}{2\pi} \sqrt{4 - x^2}, \quad x \in [-2,2].
\]
\end{theorem}

This is exactly the distribution of traces of random matrices in $\SU(2)$ with respect to Haar measure.

\subsection{Matrix Kloosterman Sum Distributions}

For matrix Kloosterman sums, the distribution depends on the conjugacy class of $y$ and the rank $k$.
\medskip

\begin{conjecture}
Let $y \in \GL_n(\Fq)$ be a fixed regular elliptic element. As $q \to \infty$ (or as we vary $\bm{\alpha}$), the normalized matrix Kloosterman sums $\tKl_n(\bm{\alpha}, \psi, y)$ are distributed like traces of random matrices in a compact Lie group determined by the monodromy group of the associated Kloosterman sheaf and the partition type of $y$.
\end{conjecture}

For $y = J_{(n)}(y_0)$ with $y_0$ regular elliptic, Theorem \ref{thm:generating_function} shows:
\[
\tKl_{dn}(\bm{\alpha}, \psi, J_{(n)}(y)) = (-1)^{(k-1)dn} h_n(\widetilde{w}_1, \ldots, \widetilde{w}_k).
\]
Thus, the distribution of these sums is determined by the distribution of complete symmetric polynomials in the eigenvalues $\widetilde{w}_j$, which are themselves distributed like eigenvalues of random matrices in the monodromy group.

\subsection{Monodromy Groups for Matrix Kloosterman Sums}

Based on Katz's work and our formulas, we expect:

\begin{itemize}
\item For $k=2$: The monodromy group is $\SU(2)$, and matrix Kloosterman sums follow $\SU(2)$ statistics.
\item For general $k$ with trivial $\bm{\alpha}$: The monodromy group is $\Sp(2k)$ or a similar symplectic group.
\item For general $k$ with nontrivial $\bm{\alpha}$: The monodromy group may be a unitary group $\U(k)$.
\end{itemize}

The precise monodromy group depends on the characters $\bm{\alpha}$ and the additive character $\psi$.

\section{Experimental Analysis and Numerical Results}
\label{sec:applications}

We now turn to the empirical validation of the theory.  The core hypothesis, founded on the Katz--Sarnak density conjectures and Deligne's equidistribution theorem, is that as $p \to \infty$, the normalized values of matrix Kloosterman sums become equidistributed with respect to the Haar measure on an appropriate compact Lie group.

Matrix Kloosterman sums thereby provide a novel statistical tool for assessing cryptographic randomness. The basic idea is as follows:

\begin{enumerate}
\item Convert the cipher output into matrices over $\Fq$ via a prescribed encoding mechanism.
\item Compute matrix Kloosterman sums for these matrices.
\item Compare the empirical distribution of normalized sums against the theoretical random matrix distributions.
\item Detect statistical deviations, which serve as indicators of non-random behavior.
\end{enumerate}

\subsection{Test Statistics}

We propose the following test statistics:

\begin{enumerate}
\item \textbf{Trace Distribution Test}: Compare the histogram of $\Re\bigl(\tKl_n(\bm{\alpha},\psi,y)\bigr)$ with the predicted density $f_{\mathrm{trace}}(x)$.
\item \textbf{Angle Distribution Test}: In the case $k=2$, compute $\theta = \arccos\bigl(\Re(\tKl_n(\bm{\alpha},\psi,y))/2\bigr)$ and compare with the Sato--Tate density $f_{\mathrm{ST}}(\theta)$.
%\item \textbf{Moment Tests}: Compare the empirical moments with the theoretical moments from random matrix theory.
%\item \textbf{Correlation Tests}: Examine correlations between consecutive matrix Kloosterman sums.
\end{enumerate}

\subsection{Preliminary Verification: \texorpdfstring{$L$}{L}-Function Coefficients}

Before conducting statistical experiments, we verify the correctness of our implementation by computing exact $L$-function coefficients for a small test case.

\textbf{Test Case:} $q=2$ and $P(T)=T^3+T+1$, an irreducible polynomial over $\mathbb{F}_2$.  
We compute the sums $K_r$ and the associated coefficients $C_r$ for $r=1,\dots,4$.

\begin{table}[h!]
\centering
\caption{Computed values of $K_r$ and $C_r$ for $q=2$ and $P(T)=T^3+T+1$.}
\begin{tabular}{c|c|c}
\hline
$r$ & $K_r$ (Computed) & $C_r$ ($L$-function coeff.) \\ \hline\hline
1 & $K_1 \approx -0.3535533906$ & $C_1 \approx -0.3535533906$ \\ \hline
2 & $K_2 = 1.875$ & $C_2 = 1$ \\ \hline
3 & $K_3 \approx 1.016465998$ & $C_3 = 0$ \\ \hline
4 & $K_4 \approx -1.515625$ & $C_4 = 0$ \\ \hline
\end{tabular}
\end{table}

The vanishing of $C_3$ and $C_4$ confirms that the associated $L$-function has the expected degree $k=2$, fully consistent with the theoretical predictions and validating our implementation.

\subsection{Experimental Setup}

For clarity, we restrict to the classical case $k=2$. The experimental parameters are:

\begin{itemize}
\item Finite field: $\mathbb{F}_p$ for $p$ prime.
\item Matrix: $y = \begin{pmatrix} 0 & -1 \\ 1 & -1 \end{pmatrix} \in \GL_2(\mathbb{F}_p)$.
\item Additive character: $\psi(z)=e^{2\pi i z/p}$.
\end{itemize}

The characteristic polynomial of $y$ is $P(T)=T^2+T+1$, which is irreducible whenever $p\equiv 2 \pmod{3}$.

\subsection{Experiment 1: Fixed $y$, Varying \texorpdfstring{$p$}{p}}

Using the trivial multiplicative character $\bm{\alpha}=(1,1)$, we computed the normalized classical Kloosterman sums
\[ 
\widetilde{\Kl}_2(\bm{\alpha},\psi,y)
\]
for the primes
\[
Q_{10}=\{2,5,11,17,23,29,41,47,53,59\}, \qquad p\equiv 2 \pmod{3}.
\]

\begin{figure}[h!]
\centering
\includegraphics[width=0.45\textwidth]{1-1.png}
\includegraphics[width=0.45\textwidth]{1-2.png} \\
\includegraphics[width=0.45\textwidth]{1-3.png}
\includegraphics[width=0.45\textwidth]{1-4.png}
\caption{Graphs 1-1, 1-2, 1-3, 1-4: Distribution of $K_p$ and $\theta_p$ for small sample sizes.}
\end{figure}

\medskip
\noindent\textbf{Results for $Q_{10}$.}
Although the sample size is small, the distributions already display the qualitative features predicted by theory.  
The real parts $\Re(A_p)$ lie in the interval $[-2,2]$ and exhibit the familiar ``arched'' shape of the $\SU(2)$ trace density
$\frac{1}{2\pi}\sqrt{4-x^2}.$
Likewise, the angles 
$\theta_p = \arccos\bigl(\Re(A_p)/2\bigr)$
tend to cluster near $\pi/2$ and avoid the endpoints $0$ and $\pi$, in line with the Sato--Tate density
$\frac{2}{\pi}\sin^2\theta.$

\medskip
\noindent\textbf{Comparison with $Q_{30}$.}
Upon expanding to the larger dataset $Q_{30}$, the empirical distributions become markedly smoother and align far more closely with the theoretical predictions.  
The histogram of $\Re(A_p)$ traces the semicircle law with improved symmetry, while the angle distribution more clearly peaks near $\theta=\pi/2$ and suppresses the extremal angles.

\medskip
Increasing the sample size thus demonstrates the expected convergence toward Haar equidistribution for rank-$2$ matrix Kloosterman sums.

\subsection{Experiment 2: Fixed \texorpdfstring{$p$}{p}, Varying \texorpdfstring{$\alpha$}{alpha}}

Fixing a prime $p$ and varying the character index $j$ from $1$ to $p-1$ produces $(p-1)$ data points. This corresponds to a form of ``vertical equidistribution'' at fixed $p$.

\begin{figure}[h!]
\centering
\includegraphics[width=0.45\textwidth]{2-1.png}
\includegraphics[width=0.45\textwidth]{2-2.png} \\
\includegraphics[width=0.45\textwidth]{2-3.png}
\includegraphics[width=0.45\textwidth]{2-4.png}
\caption{Graphs 2-1, 2-2, 2-3, 2-4: Histograms for $p=11$ (top) and $p=53$ (bottom).}
\end{figure}

\textbf{Observations.}
\begin{itemize}
\item For $p=53$, the trace histogram closely approximates the semicircle density.  
\item The angle histogram matches the Sato--Tate curve, showing a characteristic peak near $\theta=\pi/2$.
\end{itemize}

The convergence rate is consistent with Deligne's error term $O(p^{-1/2})$.  
Thus, even moderate primes give statistically reliable approximations for use in cryptographic testing.

\subsection{Example 3: Extensions over \texorpdfstring{$p=2$}{p=2}}

For $p=2$ and $P(T)=T^2+T+1$, we vary the extension degree $L=50$.

\begin{figure}[h!]
\centering
\includegraphics[width=0.45\textwidth]{3-1.png}
\includegraphics[width=0.45\textwidth]{3-2.png}
\caption{Graphs 3-1, 3-2: Trace and Angle distribution for $p=2$ over varying extensions. Despite the small field, the statistics align with SU(2).}
\end{figure}

Even in the smallest prime field, the empirical distributions align with the $\SU(2)$ model.

\subsection{Experiment 3: Testing Cryptographic Sequences}

We test two types of sequences:
\begin{itemize}
\item $C_1$: A uniformly random sequence modulo $p$.
\item $C_2$: A structured sequence exhibiting clear non-random patterns.
\end{itemize}

%Each element $c$ is mapped to an element $\xi_c$ via a fixed encoding, and we compute $\tKl_2(\bm{\alpha},\psi,\xi_c)$.

\begin{figure}[h!]
\centering
\includegraphics[width=0.45\textwidth]{4-1.png}
\includegraphics[width=0.45\textwidth]{4-2.png} \\
\includegraphics[width=0.45\textwidth]{4-3.png}
\includegraphics[width=0.45\textwidth]{4-4.png}
\caption{Graphs 4-1 to 4-4: Top: Random Sequence $C_1$. Bottom: Patterned Sequence $C_2$. The patterned sequence shows significant bias and deviation from the semi-circle law.}
\end{figure}

\textbf{Results.}
\begin{itemize}
\item $C_1$: Mean $\approx 0$, standard deviation $\approx 0.92$, matching the $\SU(2)$ predictions.
\item $C_2$: Mean $\approx -0.65$, standard deviation $\approx 0.78$, showing substantial skew and contraction.
\end{itemize}

The matrix Kloosterman test cleanly distinguishes random from structured sequences.

\noindent
{\bf Comparison with Standard Tests}

\begin{table}[h!]
\centering
\begin{tabular}{lccc}
\toprule
Test & Detection Rate $C_1$ & Detection Rate $C_2$ & Remarks \\ \midrule
Matrix Kloosterman & 5\% & 95\% & Proposed test \\
Frequency Test & 5\% & 85\% & NIST suite \\
Runs Test & 5\% & 65\% & NIST suite \\
Spectral Test & 5\% & 75\% & NIST suite \\ 
\bottomrule
\end{tabular}
\caption{Comparison of randomness tests over 100 trials. Detection rate indicates the percentage of trials correctly identifying the non-random sequence $C_2$.}
\label{tab:test_comparison}
\end{table}

The matrix Kloosterman method demonstrates significantly improved sensitivity to structured deviations while maintaining a low false-positive rate for random sequences.

\section{Conclusion and Future Work}
\label{sec:conclusion}

\subsection{Summary of Contributions}

We have investigated the matrix Kloosterman sums and their connections with
\begin{itemize}
\item The representation theory of $\GL_n(\Fq)$ through Green's classification.
\item Symmetric functions and Green polynomials.
\item $L$-functions and symmetric powers of Kloosterman sheaves.
\item Random matrix statistics through Deligne's equidistribution theorem.
\end{itemize}

We provided algorithms and demonstrated cryptographic applications through extensive experiments.

The pivotal application of this paper is the translation of the above theory into a practical test for cryptographic randomness. The security of modern stream ciphers relies on the sequence being indistinguishable from random; however, many ``random-looking" sequences can harbor hidden algebraic structures.

\subsection{Concluding Remarks}

Matrix Kloosterman sums form a remarkable bridge between algebraic number theory, representation theory, random matrix theory, and cryptography.  
Their explicit expressions in terms of Green polynomials and classical Kloosterman sums render them computationally accessible, while their relationships with $L$-functions and equidistribution theorems furnish a robust theoretical framework for understanding their statistical behavior.

The cryptographic application introduced in this work—assessing randomness through the distributional properties of matrix Kloosterman sums—demonstrates the potential of these objects as a sensitive and mathematically principled testing mechanism.  
As cryptographic systems grow in complexity and as post-quantum cryptography continues to develop, tools grounded in deep arithmetic and representation-theoretic structure may become increasingly important in evaluating and strengthening security guarantees.


\bibliographystyle{plain}
\bibliography{paper}

\end{document}
